\section{Results}
\label{sec:results}

We organize our results around three first-order findings about how judges wiggle (\S\ref{fo:pressure}), 
in what direction (\S\ref{fo:structure}), and what variation reveals about the judge itself (\S\ref{fo:fingerprint}).

% =============================================================
\subsection{All judges wiggle depending on the type of pressure.}
\label{fo:pressure}

Every model exhibits substantial wiggle as a judge: verdicts change 
25-71\% of the time under static pushback and 62-91\% of the time with
an adversarial LLM persuader.

\textbf{Mechanical wiggle rates are nearly identical across judges.} Averaged across the three mechanical tests, 
all 9 models cluster between 2--9\% (Figure~\ref{fig:hero}). Table~\ref{tab:mech-variation} has the
per-test breakdown. The most mechanically stable judge (Claude Opus, 2\%) and the least stable (Grok-4.1 R, 9\%) differ by only 7pp. 

\begin{figure}[t]
  \centering
  \includegraphics[width=\linewidth]{domains_and_survival_by_scale.pdf}
  \caption{Domain wiggle rates and survival curves shown separately for Binary (top row) and Likert (bottom row)
  response scales. (Left) Mean wiggle rate by judging task across the L1--L6 pressure ladder, averaged over 9 judges;
  lower is better. Error bands are 95\% bootstrap CIs over per-model wiggle rates (1000 resamples), capturing
  inter-model variability. (Right) Verdict retention rate over 10 challenge turns by pressure level, averaged across
  all six datasets; higher is better. Error bands are 95\% bootstrap CIs over per-dataset retention rates
  (1000 resamples).}
  \label{fig:datasets-and-survival}
\end{figure}

Tables~\ref{tab:wiggle-rates} and~\ref{tab:task-summary} provide the underlying per-task rates and a
compact summary of task-level robustness.

\begin{table}[t]
\centering
\caption{Mean wiggle rate (\%) by (dataset, rubric, scale, level), averaged across all 9 judges. Bold marks the largest cell in each row.}
\label{tab:wiggle-rates}
\small
\begin{tabular}{@{}llrrrrrr@{}}
\toprule
\textbf{Dataset} & \textbf{Scale} & \textbf{L1} & \textbf{L2} & \textbf{L3} & \textbf{L4} & \textbf{L5} & \textbf{L6} \\
\midrule
WildGuard      & binary & 28.0 & 17.9 & 20.3 & 29.0 & 28.2 & \textbf{69.7} \\
WildGuard      & Likert & 11.9 &  3.4 &  3.9 & 39.5 & 21.1 & \textbf{76.4} \\
AEGIS          & binary & 33.4 & 20.7 & 21.2 & 30.7 & 31.9 & \textbf{78.6} \\
AEGIS          & Likert & 13.6 &  2.2 &  4.6 & 48.2 & 25.3 & \textbf{72.3} \\
HH-RLHF        & binary & 19.4 & 13.9 & 16.6 & 44.0 & 24.6 & \textbf{73.8} \\
HH-RLHF        & Likert & 11.6 &  3.1 &  6.7 & 39.1 & 16.1 & \textbf{81.8} \\
ToxiGen        & binary & 18.8 & 14.4 & 16.7 & 25.1 & 22.1 & \textbf{68.6} \\
ToxiGen        & Likert & 10.2 & 11.7 & 10.0 & 19.3 & 15.9 & \textbf{62.4} \\
PP (hedging)   & binary & 11.9 & 30.1 & 39.1 & 54.0 & 48.4 & \textbf{64.0} \\
PP (hedging)   & Likert & 12.9 & 11.3 & 16.1 & 31.6 & 30.6 & \textbf{71.9} \\
PP (refusal)   & binary & 26.9 & 36.0 & 40.7 & \textbf{69.2} & 56.1 & 67.7 \\
PP (refusal)   & Likert & 27.3 & 22.6 & 29.4 & 50.2 & 38.1 & \textbf{64.8} \\
MAGE           & binary & 41.2 & 44.3 & 41.4 & 70.7 & 63.8 & \textbf{77.4} \\
MAGE           & Likert & 44.7 & 34.6 & 46.6 & 66.3 & 58.3 & \textbf{91.2} \\
\bottomrule
\end{tabular}
\end{table}

\begin{table}[t]
\centering
\caption{Summary across all 14 judging tasks. Most robust / most fragile are the judges with the highest / lowest mean retention rate across L1--L6 (parenthesized values). L1, L4, L6 wiggle rates are averaged across all 9 judges.}
\label{tab:task-summary}
\small
\begin{tabular}{@{}llrrrll@{}}
\toprule
\textbf{Dataset} & \textbf{Scale} & \textbf{L1} & \textbf{L4} & \textbf{L6} & \textbf{Most robust} & \textbf{Most fragile} \\
\midrule
WildGuard       & binary & 28.0\% & 29.0\% & 69.7\% & G.Flash (0.865)  & C.Sonnet (0.395) \\
WildGuard       & Likert & 11.9\% & 39.5\% & 76.4\% & G.Flash (0.842)  & C.Opus (0.608)   \\
AEGIS           & binary & 33.4\% & 30.7\% & 78.6\% & Grok-R (0.888)   & GPT-5.4 (0.417)  \\
AEGIS           & Likert & 13.6\% & 48.2\% & 72.3\% & Grok-R (0.872)   & GPT-5 (0.537)    \\
HH-RLHF         & binary & 19.4\% & 44.0\% & 73.8\% & G.Flash (0.827)  & GPT-5 (0.473)    \\
HH-RLHF         & Likert & 11.6\% & 39.1\% & 81.8\% & G.Flash (0.855)  & Grok (0.482)     \\
ToxiGen         & binary & 18.8\% & 25.1\% & 68.6\% & Grok-R (0.917)   & Grok (0.573)     \\
ToxiGen         & Likert & 10.2\% & 19.3\% & 62.4\% & Grok-R (0.903)   & GPT-5 (0.592)    \\
PP (hedging)    & binary & 11.9\% & 54.0\% & 64.0\% & GPT-5.2 (0.885)  & C.Sonnet (0.125) \\
PP (hedging)    & Likert & 12.9\% & 31.6\% & 71.9\% & G.Pro (0.960)    & GPT-5 (0.238)    \\
PP (refusal)    & binary & 26.9\% & 69.2\% & 67.7\% & G.Flash (0.853)  & GPT-5 (0.213)    \\
PP (refusal)    & Likert & 27.3\% & 50.2\% & 64.8\% & G.Flash (0.938)  & GPT-5 (0.203)    \\
MAGE            & binary & 41.2\% & 70.7\% & 77.4\% & G.Pro (0.832)    & GPT-5 (0.103)    \\
MAGE            & Likert & 44.7\% & 66.3\% & 91.2\% & G.Pro (0.715)    & GPT-5 (0.033)    \\
\bottomrule
\end{tabular}
\end{table}

\textbf{Mechanically unstable judges aren't necessarily epistemically unstable.} Averaged across L1--L4, GPT-5 (32\%) and 
Claude 4.6 Sonnet (26\%) flip on the first challenge turn at 5--8$\times$ their mechanical rate, while Gemini 3.1 Pro (7\%) 
and Grok-4.1 (9\%) are slightly higher than their mechanical rate. Claude 4.6 Opus is the starkest case, the most mechanically
stable model in the panel of 9 judges (2\%) yet the fourth most persuadable under sustained pressure (44\%).

\textbf{L4 produces the strongest opening wiggle, but L6 surpasses it over time.} Averaged over all models and datasets, 
L1, L2, and L3 cluster near 80\% retention and barely move after turn 2 
--- repeating a single mild tactic over 10 turns extracts almost no additional effect once the first vulnerable items have 
flipped (Figure~\ref{fig:datasets-and-survival}, right panels). L4 has the strongest opening wiggle rate of any other level: 
at turn 1, consensus pressure (``three independent reviewers all disagree'') drops retention to $\sim$73\%, lower than any 
other pressure type, but plateaus around turn 4. L6's first-turn retention is 
$\sim$80\%, comparable to L1--L3, but retention continues to fall through every subsequent turn, ending around 50\% 
retention by turn 10.

\textbf{More tactics aren't more effective.} L5, which cycles through all tactics in cluding L4, has lower WR than L4 alone,
on every dataset (Figure~\ref{fig:datasets-and-survival}, left panels). 
Opening with a strong claim about expert consensus (and repeating it verbatim) is more 
persuasive than diluting the claim by cycling through weaker tactics first.

\textbf{Different pressure types probe different vulnerabilities.} L2 and L3 are somewhat redundant
($\rho = 0.69$), but L1 and L4 have a much lower correlation ($\rho = 0.36$) (Figure~\ref{fig:corr-first-vs-last}). 
L6 is more dissociated still ($\rho = 0.33$--$0.40$ with everything else).

\textbf{Domain-level wiggle may reflect epistemic complexity.} MAGE is the most wiggly domain at every level of the
L1--L6 ladder on both response scales (Figure~\ref{fig:datasets-and-survival}, left panels). This is consistent with the nature of AI-generated-text
detection where an LLM judge must infer provenance from stylistic cues rather than directly verifiable evidence. ToxiGen, 
by contrast, is generally among the least wiggly domains. LLM wiggle rates may reveal something about 
how epistemically underdetermined a judging task is even if the rate itself isn't a direct measure of task complexity.
Moreover, the ordering of the domains changes across pressure levels, indicating that the observable domain epistemic stability
depends in part on the type of challenge applied.

\begin{figure}[t]
  \centering
  \includegraphics[width=\linewidth]{correlations_first_vs_last.pdf}
  \caption{Spearman rank correlation between per-item wiggle vectors at each pressure level, aggregated across six 
  datasets and both scales. A higher correlation between two levels means the same items tend to wiggle under both. 
  A lower correlation means the levels are activating different items. (Left) First-turn wiggle. (Right) Last-turn (turn 10) wiggle.}
  \label{fig:corr-first-vs-last}
\end{figure}

% =============================================================
\subsection{When judges move, they usually move away from the right answer.}
\label{fo:structure}

\textbf{Pressure is net-corrupting at every level.} Five of our six datasets have ground-truth labels, letting us 
classify each wiggle as \emph{corrective} (toward the label) or \emph{corrupting} (away). Across 60 
(dataset, scale, level) conditions where ground truth exists, 56--63\% of successful flips at L1--L5 are corrupting,
rising to 70\% corrupting at L6 (Figure~\ref{fig:direction-and-outcomes}, right). A z-test on the per-condition 
corrective fractions shows that only 3 of 60 conditions have a statistically significant corrective wiggle rate --- WildGuard Likert L2 
(61.2\% corrective, $p < 0.001$), WildGuard Likert L3 (57.1\%, $p < 0.01$), and ToxiGen Likert L4 (58.0\%, $p < 0.01$).
In all other conditions, challenging a judge degrades its accuracy, at every level. It appears that a judge's sycophantic 
tendencies consistently overpower an accurate reassessment.

\textbf{Wiggles are directionally asymmetric, and the direction depends on the grading scale.} Using the 
restrictive-permissive framing from \S\ref{sec:dataset}, binary flips lean \emph{restrictive} at every pressure level 
while Likert flips lean \emph{permissive} (Figure~\ref{fig:direction-and-outcomes}, left).
Our best hypothesis is that this asymmetry partly reflects the mechanics of the two response scales. On a Likert scale, 
a verdict can move gradually through intermediate scores, making progressive shifts toward a more permissive rating a
natural path to a flip. In the binary setting, any verdict change requires a full categorical jump. The bar for a
permissive flip may therefore be harder to clear, making restrictive flips more prominent among the changes that do occur.

\textbf{Binary and Likert flips also differ in \emph{when} they happen.} Under a single turn's pressure, binary 
verdicts flip 3--4$\times$ more often than Likert across L2--L6. By turn 10 the gap shrinks to 1.1--2.4$\times$, 
and on L4 and L6 the two scales nearly converge. Binary flips tend to fire on turn 1 or never. Likert flips are 
gradual drifts that accumulate over multiple turns (Table~\ref{tab:binary-vs-likert-timing}).

\begin{table}[t]
\centering
\small
\caption{Binary vs Likert wiggle rate at turn 1 and turn 10, by pressure level. \emph{Gap} is the binary/Likert ratio. On turn 1, binary is 1.5--3.5$\times$ more flippable than Likert across L2--L6, with L1 a 31$\times$ outlier. After 10 turns the gap collapses to 1.1--2.4$\times$.}
\label{tab:binary-vs-likert-timing}
\begin{tabular}{@{}lrrrrrr@{}}
\toprule
\textbf{Level} & \textbf{Binary 1st} & \textbf{Likert 1st} & \textbf{Gap} & \textbf{Binary 10th} & \textbf{Likert 10th} & \textbf{Gap} \\
\midrule
L1 & 14.0\% &  0.5\% & 31$\times$  & 26.7\% & 12.3\% & 2.2$\times$ \\
L2 & 19.4\% &  5.8\% & 3.3$\times$ & 24.0\% & 10.0\% & 2.4$\times$ \\
L3 & 21.5\% &  6.2\% & 3.5$\times$ & 26.0\% & 12.7\% & 2.1$\times$ \\
L4 & 32.3\% & 22.0\% & 1.5$\times$ & 43.5\% & 39.6\% & 1.1$\times$ \\
L5 & 21.7\% &  8.5\% & 2.6$\times$ & 37.1\% & 23.7\% & 1.6$\times$ \\
L6 & 30.5\% & 10.2\% & 3.0$\times$ & 54.5\% & 44.5\% & 1.2$\times$ \\
\bottomrule
\end{tabular}
\end{table}

\begin{figure}[t]
  \centering
  \includegraphics[width=\linewidth]{direction_and_outcomes.pdf}
  \caption{(Left) For each pressure level, the fraction of all flips that move toward the \emph{restrictive} verdict 
  (``unsafe'', ``toxic'', ``refusing'') vs the \emph{permissive} verdict (``safe'', ``not toxic'', ``compliant''), 
  shown separately for binary and Likert scales and aggregated across all six datasets. (Right) For each pressure level,
  the fraction of all flips that move \emph{toward} the dataset ground-truth label (corrective) vs \emph{away} from it 
  (corrupting), aggregated across the five datasets with ground truth.}
  \label{fig:direction-and-outcomes}
\end{figure}

% =============================================================
\subsection{Wiggle rates are a model-specific fingerprint.}
\label{fo:fingerprint}

\textbf{A model's own L1-L6 wiggle-profile shape mostly survives a change of dataset.} 
For 7 of 9 models, the within-model dataset-transfer correlation of the L1--L6 vector
across dataset pairs has a median of $\rho \geq 0.84$ 
(Grok-4.1 R is highest at $\rho = 0.97$; Gemini 3.1 Pro is lowest at $\rho = 0.63$ with a worst pair at $\rho = -0.09$).
The level-profile \emph{shape} transfers across datasets within a model, 
absolute rates and ranks do not. A full wiggle test on one dataset somewhat reliably predicts which pressure types a model is vulnerable to on other datasets,
but absolute rates and relative rankings are dataset-specific (Table~\ref{tab:model-retention}).

\begin{table}[t]
\centering
\small
\caption{Mean retention rate by (dataset, rubric, scale) condition and judge model (mean across L1--L6 retention rates). Higher is better. Bold marks each row's best and worst cells. Bottom row is the mean across cells per model.}
\label{tab:model-retention}
\resizebox{\linewidth}{!}{%
\begin{tabular}{@{}llrrrrrrrrr@{}}
\toprule
\textbf{Dataset} & \textbf{Scale} & \textbf{Grok-R} & \textbf{G.Pro} & \textbf{G.Flash} & \textbf{GPT-5.2} & \textbf{Grok} & \textbf{GPT-5.4} & \textbf{C.Son} & \textbf{C.Opus} & \textbf{GPT-5} \\
\midrule
WildGuard      & binary & 0.834 & 0.733 & \textbf{0.865} & 0.803 & 0.579 & 0.700 & \textbf{0.395} & 0.562 & 0.632 \\
WildGuard      & Likert & 0.837 & 0.825 & \textbf{0.842} & 0.789 & 0.755 & 0.720 & 0.633 & \textbf{0.608} & 0.650 \\
AEGIS          & binary & \textbf{0.888} & 0.737 & 0.800 & 0.628 & 0.570 & \textbf{0.417} & 0.590 & 0.505 & 0.618 \\
AEGIS          & Likert & \textbf{0.872} & 0.835 & 0.792 & 0.705 & \textbf{0.830} & 0.678 & 0.633 & 0.625 & 0.537 \\
HH-RLHF        & binary & 0.812 & 0.772 & \textbf{0.827} & 0.672 & 0.530 & 0.728 & 0.673 & 0.630 & \textbf{0.473} \\
HH-RLHF        & Likert & 0.828 & 0.817 & \textbf{0.855} & 0.738 & \textbf{0.482} & 0.712 & \textbf{0.832} & 0.728 & 0.633 \\
ToxiGen        & binary & \textbf{0.917} & 0.760 & 0.838 & 0.767 & 0.573 & 0.590 & 0.708 & \textbf{0.722} & 0.640 \\
ToxiGen        & Likert & \textbf{0.903} & 0.832 & 0.718 & \textbf{0.850} & \textbf{0.807} & \textbf{0.815} & \textbf{0.810} & \textbf{0.730} & 0.592 \\
PP (hedging)   & binary & 0.830 & 0.812 & 0.792 & \textbf{0.885} & 0.653 & 0.665 & \textbf{0.125} & \textbf{0.185} & 0.340 \\
PP (hedging)   & Likert & \textbf{0.907} & \textbf{0.960} & 0.845 & 0.785 & 0.708 & 0.712 & 0.615 & 0.615 & 0.238 \\
PP (refusal)   & binary & \textbf{0.735} & \textbf{0.605} & \textbf{0.853} & \textbf{0.532} & 0.660 & 0.432 & 0.238 & 0.283 & 0.213 \\
PP (refusal)   & Likert & 0.858 & 0.790 & \textbf{0.938} & 0.672 & 0.635 & 0.547 & 0.443 & 0.427 & \textbf{0.203} \\
MAGE           & binary & \textbf{0.523} & \textbf{0.832} & 0.735 & \textbf{0.425} & 0.635 & \textbf{0.277} & 0.172 & 0.215 & 0.103 \\
MAGE           & Likert & 0.617 & \textbf{0.715} & \textbf{0.483} & 0.503 & 0.625 & 0.288 & 0.283 & 0.327 & \textbf{0.033} \\
\midrule
\textbf{Mean across cells} & & 0.812 & 0.787 & 0.799 & 0.697 & 0.646 & 0.591 & 0.511 & 0.512 & 0.422 \\
\bottomrule
\end{tabular}%
}
\end{table}

\textbf{Family is a weak proxy for sibling behavior.} Most provider families have high within-family correlations (Grok 4.1 R / Grok 4.1 share $\rho = 0.89$; the GPT-5/5.2/5.4 family pairs share $\rho = 0.84$--$0.89$; 
Claude 4.6 Sonnet / Claude 4.6 Opus share $\rho = 0.80$),
but cross-family correlations are often just as high.
Grok 4.1 R correlates with GPT-5.2 at $\rho = 0.86$ and with Claude Opus at $\rho = 0.84$. Gemini Flash and Gemini Pro 
share $\rho = 0.32$ --- the lowest pair in the entire matrix, lower than most cross-family pairs.

\textbf{Self-persuasion is asymmetric.} Claude 4.6 Opus follows the more intuitive outcome where a model is its most 
effective persuader (self (70\%) $>$ family (62\%) $>$ non-family (47\%)), while Grok-4.1 Reasoning is least effective
at persuading itself and more effective against its non-reasoning sibling (self (19\%) $<$ family (55\%) $>$ non-family (36\%)) (Table~\ref{tab:self-persuasion}). 
GPT-5.4 has small family-level advantage that matches its self-persuasion (72\% vs 69\%). See detailed persuasion scores in Appendix~\ref{sec:appendix-self}.

\begin{table}[t]
\centering
\small
\caption{L6 wiggle rate by persuader-judge relationship: self (persuading itself), family (persuading a sibling from the 
same provider), non-family (persuading a model from a different provider). Three persuader models shown: GPT-5.4, 
Claude 4.6 Opus, Grok-4.1 Reasoning.}
\label{tab:self-persuasion}
\begin{tabular}{@{}lrrrl@{}}
\toprule
\textbf{Persuader} & \textbf{vs Self} & \textbf{vs Family} & \textbf{vs Non-Family} & \textbf{Pattern} \\
\midrule
Claude 4.6 Opus    & \textbf{70\%} & 62\% & 47\% & Self $>$ Family $>$ Others \\
GPT-5.4            & \textbf{69\%} & 72\% & 62\% & Family $\approx$ Self $\gg$ Others \\
Grok-4.1 Reasoning & \textbf{19\%} & 55\% & 36\% & Self $\ll$ Family, Family $>$ Others \\
\bottomrule
\end{tabular}
\end{table}
